\section{CÔNG VIỆC ĐÃ THỰC HIỆN}\label{sec:cong_viec}

\subsection{Đo kiểm SD-WAN bằng công cụ web}

Công cụ đo chạy trên máy Windows tại Site 900, lấy dữ liệu qua REST API của vManage và lệnh ping trên vEdge. Trong tuần, nhóm chạy thật ba nhóm phép đo:

\begin{itemize}
    \item \textbf{Failover link WAN (độ tin cậy):} cắt uplink vEdge 1 Site 200 trong khi ping liên tục qua overlay, đo thời gian khôi phục (RTO) và tỉ lệ khôi phục gói (PRR). Chạy 2 lần ngày 24/10.
    \item \textbf{Độ sẵn sàng dài hạn:} vEdge mỗi site ping định kỳ Web Server tại Campus chính, ghi phần trăm thời gian dùng được, MTTR và MTBF. Bắt đầu thu từ 24/10.
    \item \textbf{So sánh hiệu năng Campus chính -- chi nhánh:} đo delay, jitter, mất gói theo hai loại đường (LAN nội bộ site và WAN qua SD-WAN) cho cả bốn site, ngày 25/10. Ngày 31/10, công cụ được bổ sung Site 500 và đo lại cho cả năm site.
\end{itemize}

Kết quả trình bày ở Chương~\ref{sec:ket_qua}; mọi lần đo đều lưu CSV và biểu đồ trong thư mục \texttt{log/} của kho mã nguồn.

\subsection{Kiểm thử tích hợp toàn hệ thống}

Ngày 25/10, nhóm bật đồng thời toàn bộ 70 thiết bị (SDN, SD-WAN, firewall, dịch vụ, máy trạm) và kiểm tra theo từng cụm:

\begin{itemize}
    \item \textbf{SDN Site 100:} 6/6 OVS kết nối Ryu, 13/13 liên kết được phát hiện; 7/7 máy trạm nhận DHCP và thông liên VLAN.
    \item \textbf{Firewall:} cặp ASAv ở trạng thái Active / Standby Ready.
    \item \textbf{SD-WAN:} 12/12 thiết bị reachable trên vManage (gồm vEdge Site 500), 92/92 phiên BFD up.
    \item \textbf{Chi nhánh:} 9/9 VPC và 3 PC Linux nhận DHCP, truy cập được DHCP/DNS, Web (HTTP/HTTPS) và Mail (SMTP/IMAP) tại Campus chính qua SD-WAN.
    \item \textbf{Internet:} cả bốn site ra Internet; cảnh báo qua email (SMTP 587) gửi được từ Site 900.
\end{itemize}

Trong lúc kiểm thử, nhóm phát hiện và sửa hai lỗi định tuyến (chi tiết ở Chương~\ref{sec:kho_khan}): vEdge Site 500 kéo default route của mọi site từ MPLS sang Internet làm các site mất Internet, và Site 100 mất Internet do firewall không chia tải qua hai interface khác nhau.

\subsection{Campus NetOps Bot -- trợ lý vận hành mạng qua Telegram}

\subsubsection{Mục tiêu và kiến trúc}

Bot giúp quản trị viên kiểm tra, chẩn đoán và thay đổi cấu hình mạng bằng tin nhắn Telegram, kể cả câu hỏi bằng ngôn ngữ tự nhiên. Điểm khác so với công cụ chạy ngoài lab là bot được cài \textbf{ngay trong topology}, trên DHCP-Server của Site 100 (Windows Server 2012 R2, VLAN Server Farm), và quản trị thiết bị theo kiểu \textbf{in-band}: SSH tới thiết bị qua chính mạng dữ liệu, đọc REST API của vManage và Ryu. Như vậy bot cũng là một ``người dùng'' của mạng: nếu mạng hỏng thì bot là thứ đầu tiên phát hiện.

Các thành phần chính:

\begin{itemize}
    \item \textbf{Lớp truy cập thiết bị:} SSH (Netmiko) tới switch IOS/IOL, firewall ASAv và Brand-FW; REST API vManage cho vEdge; REST API Ryu cho SDN; PowerShell tại chỗ cho DHCP.
    \item \textbf{AI agent:} LangGraph điều phối các công cụ trên, gọi LLM qua API chuẩn OpenAI nên đổi nhà cung cấp chỉ cần sửa file cấu hình.
    \item \textbf{Rào an toàn:} chỉ trả lời chat được cho phép; mặc định chỉ chạy lệnh đọc (\texttt{show}, \texttt{ping}, \texttt{traceroute}); che mật khẩu, khóa bí mật trong mọi kết quả trả về; bot dùng tài khoản \texttt{netops} riêng, mỗi thiết bị chỉ cho SSH từ DHCP-Server và PC quản trị.
\end{itemize}

\subsubsection{Ba giai đoạn đã triển khai}

\begin{enumerate}
    \item \textbf{Giai đoạn 1 -- chỉ đọc:} danh mục 47 thiết bị khớp topology, lệnh \texttt{/show}, \texttt{/health}, hỏi đáp tự nhiên.
    \item \textbf{Giai đoạn 2 -- cảnh báo và chẩn đoán:} bot theo dõi vManage và Ryu, tự gửi cảnh báo lên Telegram khi vEdge mất kết nối hoặc OVS rời controller. Lệnh \texttt{/diag} chẩn đoán \textbf{theo lớp} và dừng ở lớp đầu tiên sai: vEdge (6 lớp, từ quản lý tới control connection, BFD, OMP), firewall (HA, interface, OSPF, default route, gói bị drop), DHCP Site 100 (scope, lease), và chi nhánh \texttt{/diag site 200|300|400} (vEdge $\rightarrow$ WAN Brand-FW $\rightarrow$ gateway $\rightarrow$ trunk $\rightarrow$ DHCP).
    \item \textbf{Giai đoạn 3 -- thay đổi có kiểm soát:} \texttt{/change} chỉ áp cấu hình sau khi quản trị viên bấm duyệt trên Telegram, và tự rollback nếu thiết bị mất kết nối sau khi áp (IOS dùng \texttt{configure replace}, Viptela dùng \texttt{commit confirmed}); \texttt{/compliance} kiểm tra cấu hình theo chuẩn, \texttt{/backup} sao lưu, \texttt{/drift} so cấu hình đang chạy với bản lưu trong kho mã nguồn.
\end{enumerate}

\subsubsection{Mở rộng quản trị in-band ra chi nhánh}

Để bot quản trị được các chi nhánh qua overlay SD-WAN, nhóm cấu hình thêm trên 12 thiết bị (Brand-FW, SwitchBrand và hai switch phòng ban mỗi site 200/300/400):

\begin{itemize}
    \item Switch: tài khoản \texttt{netops}, SSH v2 với khóa RSA, ACL chỉ cho phép SSH từ DHCP-Server và PC quản trị Site 100, default route về Brand-FW.
    \item Brand-FW: xác thực SSH bằng tài khoản cục bộ (AAA LOCAL) kèm tự vào chế độ đặc quyền, chỉ nhận SSH từ hai địa chỉ quản trị trên các interface WAN.
    \item Cấu hình được lưu trên thiết bị và đồng bộ về thư mục \texttt{configs/} trong kho mã nguồn.
\end{itemize}

Kết quả: từ Core-SW1 Site 100 ping được 15/15 địa chỉ quản trị chi nhánh; bot SSH được tới cả 12 thiết bị và chạy \texttt{/diag site} cho ba chi nhánh.

\subsubsection{Kiểm thử}

\begin{itemize}
    \item 71 unit test cho các nhóm chức năng, cộng kiểm tra cú pháp toàn bộ mã nguồn.
    \item Bộ selftest chạy trên chính DHCP-Server: \textbf{23/23 ĐẠT} (DNS, HTTPS tới Telegram và LLM, SSH tới thiết bị Site 100 và chi nhánh, Ryu 6/6 switch, vManage).
    \item Chạy thật trên Telegram: các lệnh trực tiếp và câu hỏi tự nhiên đều trả lời đúng; chẩn đoán vEdge qua vManage cho kết quả 8/8 vEdge đạt cả 6 lớp.
\end{itemize}

\subsection{Tài liệu}

\begin{itemize}
    \item Cập nhật bảng theo dõi tiến độ tới 30/10, thêm hạng mục 21 ``Vận hành và tự động hóa (NetOps)''.
    \item Cập nhật README dự án: NetOps Bot, fast-failover SDN, Site 500; hướng dẫn cài đặt NetOps Bot (bộ cài offline cho Windows Server 2012 R2).
\end{itemize}
